% Fichier complété par tools/complete_missing_corrections.py.
% Source : RDM/RDM-01-Cohesion/529_Passerelle/529_Passerelle.tex
% Statut : rédigé (7 question(s)).
\begin{corrigebox}[Corrigé rédigé]
\CorrigeQuestion{1}
On choisit l'axe $x$ le long de chaque poutre, l'origine aux liaisons et
            les coordonnées des points d'application. Le câble est paramétré par sa
            longueur $L_c$, sa section $S_c$ et son module $E_c$.
\CorrigeQuestion{2}
On isole l'ensemble de la passerelle. Les équations
            $\sum\vec F=\vec0$ et $\sum\vec M_A=\vec0$ donnent la tension du câble et les
            réactions en A; les actions réciproques entre poutres sont obtenues par
            isolement successif.
\CorrigeQuestion{3}
Une coupure à l'abscisse $x$ dans chaque poutre donne
            $\{T_{coh}\}=-\sum\{T_{ext}\}$; les expressions sont établies séparément de
            part et d'autre des charges et des nœuds.
\CorrigeQuestion{4}
On trace $N(x)$, $T(x)$ et $M_f(x)$ en respectant les valeurs aux
            extrémités, les sauts dus aux forces concentrées et les pentes imposées par les
            charges réparties.
\CorrigeQuestion{5}
L'allongement du câble est
            $\boxed{\Delta L_c=TL_c/(E_cS_c)}$ pour une tension uniforme.
\CorrigeQuestion{6}
Pour une petite rotation $\theta_A$, le déplacement de B est
            $\delta_B\simeq AB\,\theta_A$ perpendiculairement à la poutre. Le déplacement
            de C s'obtient par composition des petites rotations et des allongements des
            deux tronçons.
\CorrigeQuestion{7}
Pour une section composée,
            $I_{Oy}=\sum(I_{G_i y}+S_i z_i^2)$ et
            $I_{Oz}=\sum(I_{G_i z}+S_i y_i^2)$ (Huygens). Les évidements sont soustraits.
\end{corrigebox}
